\color{unidarkgreen}{\subsection{Контроль цепей напряжения}}
\color{black}

\begin{enumerate}[label=\arabic{section}.\arabic{subsection}.\arabic{enumi}, labelsep=4pt, leftmargin=0pt, itemindent=57pt, itemsep=0pt, parsep=5pt] % если что можно крутить itemsep


%%%%%%%%%%%%%%%%%%%%%%%%%%%%%%%%%%%%%%%%%%%%%%%%%% Контроль цепей напряжения %%%%%%%%%%%%%%%%%%%%%%%%%%%%%%%%%%%%%%%%%%%%%%%%%%
\item Общие данные

\begin{enumerate}[label=\arabic{section}.\arabic{subsection}.\arabic{enumi}.\arabic*, labelsep=4pt, leftmargin=0em, itemindent=65pt, parsep=0pt]

\item
Логика контроля цепей напряжения предназначена:
\begin{itemize}
\item для диагностики неисправности во вторичных цепях ТН;
\item для контроля отсутствия напряжения на СШ для ввода очувствления ДЗШ при опробовании;
\item для контроля наличия напряжения на СШ для сброса автоматического очувствления ДЗШ при успешном АПВ шин;
\item для запрета АПВ при наличии напряжения на СШ (при неполнофазном или полнофазном отказе выключателя питающего присоединения).
\end{itemize}

\medskip

\begin{figure}[H]
\centering
\includegraphics[width=0.9\textwidth,height=0.9\textheight,keepaspectratio]{\fbpath/661.1731 КЦН/img/KCN.pdf}
\captionof{figure}{Функциональный блок КЦН 1С}\label{pic:KCN}
\end{figure}

\medskip

\item
Ввод и вывод в работу контроля исправности цепей напряжения для каждого ТН на СШ осуществляется либо посредством программного ключа \textbf{«КЦН 1С»} (\textbf{«КЦН 2С»}), либо с использованием внешнего сигнала \textit{«Вывод КЦН 1С»} (\textit{«Вывод КЦН 2С»}).

\item
При снижении значения хотя бы одного линейного напряжения ниже уставки \textbf{«ИО отсутствия U 1C»} (\textbf{«ИО отсутствия U 2C»}), либо при срабатыванием ИО наличия напряжения ОП (\textbf{«ИО наличия ОП U 1C»}, \textbf{«ИО наличия ОП U 2C»}) через выдержку времени \textbf{«Ткцн 1С»} (\textbf{«Ткцн 2С»}) формируется сигнал \textit{«Неиспр. ЦН 1С»} (\textit{«Неиспр. ЦН 2С»}).
Данный сигнал свидетельствует о наличии неисправности цепей напряжения ТН на СШ.

\item
Сигнал \textit{«Отсутствие U 1С»} (\textit{«Отсутствие U 2С»}) формируется, если значения всех линейных напряжений опускается ниже уставки \textbf{«ИО отсутствия U 1C»} (\textbf{«ИО отсутствия U 2C»}) и одновременно отсутствует срабатывание ИО наличия напряжения ОП.
Данный сигнал свидетельствует об отсутствии напряжения и переходного режима на СШ.

\item
Сигнал \textit{«Наличие U при откл. 1С»} (\textit{«Наличие U при откл. 2С»}) формируется, если значение хотя бы одного линейного напряжения превысят уставку \textbf{«ИО наличия U 1C»} (\textbf{«ИО наличия U 2C»}), либо при срабатывании ИО наличия напряжения ОП (\textbf{«ИО наличия ОП U 1C»}, \textbf{«ИО наличия ОП U 2C»}).
Данный сигнал свидетельствует о наличии напряжения на отключенной СШ и используется в логике запрета АПВ для факта наличия неполнофазного (полнофазного) отказа выключателя.

\item
Сигнал \textit{«Наличие U 1С»} (\textit{«Наличие U 2С»}) формируется, если значения всех линейных напряжений превысят уставку \textbf{«ИО наличия U 1C»} (\textbf{«ИО наличия U 2C»}) и одновременно отсутствует срабатывание ИО наличия напряжения ОП.
Данный сигнал предназначен для сброса автоматического очувствления после успешного АПВ первого присоединения и свидетельствует о нормальном режиме работы и отсутствии переходного режима на СШ.

\item
Уставки функции КЦН представлены в таблице \ref{desc:tbl1}.

\small
\begin{longtable}{|>{\centering\arraybackslash}m{5.3cm}|>{\centering\arraybackslash}m{3.3cm}|>{\centering\arraybackslash}m{4.2cm}|>{\centering\arraybackslash}m{1.8cm}|>{\centering\arraybackslash}m{1cm}|}
% Первая страница: полный заголовок
\caption{Перечень уставок функции КЦН\hfill\vspace{-0.5\baselineskip}}\label{desc:tbl1}\\
\hline
\rowcolor{gray!30}
Параметр & Условное обозначение на схеме & Значение/ Диапазон & Единица измерения & Шаг \\
\hline
\endfirsthead

% Промежуточные страницы: «Продолжение таблицы»
\caption*{\hspace{3pt}\emph{Продолжение таблицы \ref{desc:tbl1}}\hfill\ }\vspace{-0.5\baselineskip} \\
\hline
\rowcolor{gray!30}
Параметр & Условное обозначение на схеме & Значение/ Диапазон & Единица измерения & Шаг \\
\hline
\endhead

%===m>ДЗШ|-
%===m
\end{longtable}
\normalsize	

\end{enumerate}

\end{enumerate}
